% Generated draft: requires mathematical and source review.
\documentclass[11pt]{article}
\usepackage[margin=25mm]{geometry}
\usepackage{fontspec}
\setmainfont{Latin Modern Roman}
\usepackage{amsmath,amssymb,mathtools,unicode-math}
\setmathfont{Latin Modern Math}
\usepackage{enumitem,needspace,xurl,hyperref}
\setlength{\parindent}{0pt}
\setlength{\parskip}{5pt}
\setlength{\emergencystretch}{3em}
\newcommand{\N}{\mathbb{N}}
\newcommand{\Z}{\mathbb{Z}}
\newcommand{\Q}{\mathbb{Q}}
\newcommand{\R}{\mathbb{R}}
\newcommand{\C}{\mathbb{C}}
\newcommand{\F}{\mathbb{F}}
\newcommand{\A}{\mathbb{A}}
\newcommand{\PP}{\mathbb{P}}
\newcommand{\E}{\mathbb{E}}
\newcommand{\eps}{\varepsilon}
\newcommand{\dd}{\,\mathrm{d}}
\begin{document}
\section*{UnsolvedMath ID: 30003565}
\textbf{Source problem code:} catalogue entry
\subsection{Definitions and mathematical statement}
Let \(K\subset S^3\) be a hyperbolic knot, meaning that its complement \(M_K=S^3\setminus K\) admits a complete hyperbolic metric of sectional curvature \(-1\) and finite volume. Write \(\operatorname{Vol}(M_K)\) for that hyperbolic volume. For a positive integer \(N\), put \(q_N=e^{2\pi i/N}\), and let \(J_N(K;q_N)\) denote the root-of-unity specialization of the \(N\)-colored Jones invariant named in the supplied record. Its normalization, framing convention, and precise indexing convention are not specified there. Write \(\operatorname{CS}(M_K)\) for the Chern--Simons quantity appearing in the record; its normalization, orientation convention, period, and choice of representative are likewise unspecified.

The supplied complexified volume conjecture asks whether, for every hyperbolic knot \(K\subset S^3\),
\[
\lim_{N\to\infty}\frac{2\pi}{N}\log J_N\!\left(K;e^{2\pi i/N}\right)
\stackrel{?}{=}
\operatorname{Vol}(M_K)+i\operatorname{CS}(M_K).
\]
Here \(N\) tends to infinity through positive integers. This display records the intended target, but the supplied data do not determine a well-defined equality of complex numbers. To interpret its left-hand side, one must specify logarithm values \(L_N\in\C\) satisfying
\[
 e^{L_N}=J_N(K;q_N)
\]
for all sufficiently large \(N\), which requires eventual nonvanishing, and specify how these logarithm values are selected across \(N\). Changing \(L_N\) by \(2\pi i k_N\), with \(k_N\in\Z\), changes the scaled term by \(4\pi^2 i k_N/N\), so the limiting imaginary part is not determined without a selection rule. The record supplies neither such a rule nor an eventual nonvanishing qualification. It also does not specify whether the right-hand side uses a chosen lift of the Chern--Simons quantity or whether the proposed equality is intended modulo a period. These unresolved choices must be fixed before the displayed target becomes a precise conjecture; no choices are imposed in this draft.
\subsection{Short English statement}
For every hyperbolic knot, the supplied conjecture relates the asymptotic logarithm of its colored Jones invariant at \(e^{2\pi i/N}\) to the hyperbolic volume and Chern--Simons quantity of its complement. The logarithm selection, invariant normalizations, and meaning of the Chern--Simons term remain unspecified, so the complex equality requires clarification.
\subsection{Sources}
\begin{itemize}
\item Local UnsolvedMath registry, permanent ID \(30003565\), source problem code \texttt{catalogue entry}, title: \emph{Complexified Volume Conjecture for Hyperbolic Knots}. The supplied record belongs to the Oberwolfach Reports Open Problems collection and lists proposed year 2017; no individual proposer is supplied.
\item Supplied report citation: \emph{Low-dimensional Topology and Number Theory} (2018), Oberwolfach Reports, pp.\ 2363--2426, DOI \texttt{10.4171/owr/2017/38}. \url{https://doi.org/10.4171/owr/2017/38}. The supplied verification note says that the official report was checked and that the reconstruction restored the root-of-unity evaluation and identified the geometric quantities with invariants of the knot complement. This is a source-record assertion, not an independent verification by this draft.
\item Supplied literature source: \emph{The colored Jones polynomial and the simplicial volume}. \url{https://doi.org/10.1007/BF02392716}.
\item Supplied literature source: \emph{Low dimensional Topology and Number Theory}. \url{https://doi.org/10.4171/owr/2017/38}.
\item Literature assessment dated exactly \texttt{2026-08-21}: supplied dated curation, not independently verified. The record labels the problem open and describes the truncated record as identifying the general quantum knot invariant volume relation, whose general hyperbolic knot case it assesses as unresolved. Its background explicitly requires independent verification of time-sensitive status and cited claims. This supplied assessment does not settle the formulation issues in the complexified equality.
\end{itemize}
\end{document}
